\begin{table}[t]\centering\small
\caption{The nine Pacific Rim earthquakes, their study domains and landslide statistics on the 30~m grid. Prevalence is the percentage of domain cells that contain a landslide point. PGA values are from the USGS ShakeMap used as model input.}\label{tab:events}
\resizebox{\textwidth}{!}{\begin{tabular}{lrrrrrrr}\toprule
Event & M\textsubscript{w} & Domain (km$^2$) & Landslides & Positive cells & Prevalence (\%) & PGA$_{\max}$ (g) & mean PGA at landslides (g)\\\midrule
Kaik\={o}ura 2016 (NZ) & 7.8 & 10\,832 & 29\,519 & 28\,474 & 0.237 & 1.61 & 0.64 \\
Northridge 1994 (USA) & 6.7 & 4\,377 & 11\,111 & 10\,925 & 0.225 & 0.87 & 0.48 \\
Niigata-Chuetsu 2004 (JP) & 6.6 & 410 & 4\,862 & 4\,860 & 1.067 & 0.82 & 0.70 \\
Iwate-Miyagi 2008 (JP) & 6.9 & 777 & 4\,164 & 4\,093 & 0.474 & 1.20 & 0.96 \\
Ais\'{e}n 2007 (CL) & 6.2 & 1\,247 & 517 & 512 & 0.037 & 0.12 & 0.06 \\
Denali 2002 (USA) & 7.9 & 15\,748 & 1\,579 & 1\,579 & 0.009 & 0.74 & 0.49 \\
Cinchona 2009 (CR) & 6.1 & 637 & 4\,826 & 4\,743 & 0.670 & 0.27 & 0.19 \\
Porgera 2018 (PNG) & 7.5 & 24\,306 & 10\,469 & 10\,455 & 0.039 & 1.05 & 0.53 \\
Palu 2018 (ID) & 7.5 & 3\,665 & 7\,063 & 7\,007 & 0.172 & 0.88 & 0.43 \\
\bottomrule\end{tabular}}\end{table}
